% Part: incompleteness
% Chapter: theories-computability
% Section: oconsis-ext-of-q-undec

\documentclass[../../../include/open-logic-section]{subfiles}

\begin{document}

\olfileid{inc}{tcp}{oqn}

\olsection{$\omega$-بے تضاد توسیعاں جیہڑیاں $\Th{Q}$ نوں شامل کردیاں نیں، ناقابلِ فیصلہ نیں}

\begin{explain}
$\Th{Q}$ دی قابلِ حساب گنتی وار فہرست بنن جوگ سیٹاں وچ مکمل ہون
دا ثبوت ایس گل اُتے ٹکیا سی کہ $\Th{Q}$ وچ ہر ثابت ہو سکّن والا
جملہ قدرتی عدداں بارے “سچ” اے۔ اگلی تعریف تے قضیہ سچائی دے اوہناں خاص
پہلوآں دی نشان دہی کردے نیں جیہڑے پچھلے ثبوت لئی لوڑیںدے سن۔
ایس طرح نتیجہ ہور مضبوط ہو جاندا اے۔ ایس قضیے اُتے بہتی دیر
نہ لاؤ، کیوں جے اسیں چھیتی ایہنوں ہور مضبوط کراں گے۔ ایہنوں
شامل کرن دا خاص سبب تاریخی اے۔ گوڈل نے اپنے اصل مقالے وچ
$\omega$-بے تضادی دا تصور ورتیا سی۔ تھوڑے چِر پچھوں
$\omega$-بے تضادی دی جگہ عام بے تضادی رکھ کے اوہدا نتیجہ
ہور مضبوط کر دتا گیا۔
\end{explain}

\begin{defn}
\ollabel{thm:oconsis-q}
نظریہ $\Th{T}$، $\omega$-بے تضاد اے جے ہِٹھلی شرط پوری ہووے۔
جے $\lexists[x][!A(x)]$ کوئی وی جملہ اے تے $\Th{T}$،
$\lnot
!A(\num 0)$، $\lnot !A(\num 1)$، $\lnot !A(\num 2)$، \dots
سبھ ثابت کردا اے، تاں $\Th{T}$، جملہ
$\lexists[x][!A(x)]$ ثابت نہیں کردا۔
\end{defn}

\begin{thm}
فرض کرو $\Th{T}$ اک $\omega$-بے تضاد نظریہ اے جیہڑا $\Th{Q}$
نوں شامل کردا اے۔ تاں $\Th{T}$ قابلِ فیصلہ نہیں۔
\end{thm}

\begin{proof}
جے $\Th{T}$ وچ $\Th{Q}$ شامل اے، تاں $\Th{T}$ قابلِ حساب
فنکشناں تے تعلقاں دی نمائندگی کردا اے۔ سانوں پچھلے ثبوت وچ
بس تھوڑی جیہی تبدیلی کرنی پئے گی۔ پہلاں وانگ جے $x \in K$،
تاں $\Th{T}$، جملہ $\lexists[s][!A_T(\num
  x,\num x, s)]$ ثابت کردا اے۔ دوجے پاسے فرض کرو $\Th{T}$
جملہ $\lexists[s][!A_T(\num x, \num x, s)]$ ثابت کردا اے۔
فیر $x$ سیٹ $K$ وچ ضرور اے۔ جے ایہ اوہدے وچ نہ ہووے، تاں
مشین $x$ دی انپٹ $x$ اُتے کوئی رُکّن والی حسابی کارروائی نہیں۔
کیوں جے $!A_T$، کلینی دے تعلق $T$ دی نمائندگی کردا اے،
ایس لئی $\Th{T}$، فارمولے $\lnot !A_T(\num x, \num x, \num 0)$،
$\lnot !A_T(\num x, \num
x, \num 1)$، \dots سبھ ثابت کردا اے۔ ایس طرح $\Th{T}$،
$\omega$-بے تضاد نہیں رہندا۔
\end{proof}

\end{document}
